\documentclass[12pt,a4paper]{article}

\usepackage[utf8]{inputenc}
\usepackage[T1]{fontenc}
\usepackage[brazil]{babel}
\usepackage[margin=2.5cm]{geometry}
\usepackage{graphicx}
\usepackage{hyperref}
\usepackage{xcolor}
\usepackage{booktabs}
\usepackage{enumitem}
\usepackage{titlesec}
\usepackage{parskip}
\usepackage{listings}

\hypersetup{
    colorlinks=true,
    linkcolor=blue!60!black,
    urlcolor=blue!60!black,
    citecolor=blue!60!black
}

\titleformat{\section}{\normalfont\Large\bfseries}{\thesection.}{0.6em}{}
\titleformat{\subsection}{\normalfont\large\bfseries}{\thesubsection.}{0.5em}{}

\begin{document}

\begin{titlepage}
    \begin{center}
        \includegraphics[width=0.22\textwidth]{ufsc_logo.png}
        \vspace{1cm}

        {\large \textbf{UNIVERSIDADE FEDERAL DE SANTA CATARINA}}\\
        {\large Departamento de Engenharia El\'etrica e Eletr\^onica}\\
        {\large Projeto N\'ivel I em Cont.\ e Proc.\ de Sinais I}

        \vspace{3cm}

        {\LARGE \textbf{Proposta de Projeto}}\\
        \vspace{0.3cm}
        {\large Sistema Web de Reconhecimento Facial com Prova de Vida}

        \vspace{1cm}

        {\Large \textbf{FaceID Lab}}

        \vfill

        \textbf{Acad\^emico:}\\
        Mateus Gripp Pereira\\[0.4cm]
        \textbf{Professor:} Joceli Mayer\\
        \textbf{Disciplina:} EEL7815 -- Semestre 2026.1

        \vfill

        Florian\'opolis\\
        Mar\c{c}o de 2026
    \end{center}
\end{titlepage}

\tableofcontents
\newpage

% ==========================================================================
\section{T\'itulo e Equipe}

\textbf{T\'itulo:} FaceID Lab -- Sistema Web de Cadastro e Identifica\c{c}\~ao Facial com Prova de Vida Baseado em Descritores Profundos.

\subsection*{Equipe e Responsabilidades}

\begin{itemize}[leftmargin=*]
    \item \textbf{Mateus Gripp Pereira} (\textit{respons\'avel \'unico}): concep\c{c}\~ao,
    arquitetura e implementa\c{c}\~ao integral do sistema, incluindo:
    \begin{itemize}
        \item \emph{Backend} Node.js + Express + Prisma/SQLite, com integra\c{c}\~ao
        do \texttt{@vladmandic/face-api} (detec\c{c}\~ao, \emph{landmarks} e descritor
        de 128-D) e API REST cobrindo cadastro com prova de vida, cadastro simples
        (legado), identifica\c{c}\~ao, listagem e consulta individual;
        \item \emph{Frontend} React + TypeScript + Vite, com captura via webcam
        (\texttt{getUserMedia}) e telas de cadastro, identifica\c{c}\~ao e listagem;
        \item Pipeline de \emph{prova de vida} ativa (an\'alise temporal de \emph{yaw},
        EAR e consist\^encia de descritores entre \emph{frames});
        \item Defesas de servidor: verifica\c{c}\~ao \emph{anti-adultera\c{c}\~ao}
        (re-extra\c{c}\~ao do descritor no servidor e compara\c{c}\~ao com o reportado pelo
        cliente) e \emph{detec\c{c}\~ao de cadastro duplicado} contra a base existente;
        \item Suporte ao upload opcional de imagem de \emph{documento de identidade}
        para auditoria;
        \item Constru\c{c}\~ao da base de testes interna, calibra\c{c}\~ao do limiar de
        decis\~ao e avalia\c{c}\~ao experimental (FAR/FRR);
        \item Reda\c{c}\~ao do relat\'orio final e apresenta\c{c}\~ao.
    \end{itemize}
\end{itemize}

% ==========================================================================
\section{Introdu\c{c}\~ao}

O reconhecimento facial \'e uma das aplica\c{c}\~oes mais difundidas de \emph{Processamento Digital
de Imagens} (PDI) e \emph{Vis\~ao Computacional}, presente em controle de acesso, autentica\c{c}\~ao
banc\'aria, registro de presen\c{c}a, atendimento personalizado e seguran\c{c}a p\'ublica. O problema
combina v\'arias etapas cl\'assicas estudadas na disciplina: aquisi\c{c}\~ao e amostragem da imagem,
filtragem e normaliza\c{c}\~ao, detec\c{c}\~ao de regi\~oes de interesse (faces), extra\c{c}\~ao de
descritores de caracter\'isticas e classifica\c{c}\~ao por dist\^ancia em um espa\c{c}o de
\emph{embeddings}.

Apesar de o reconhecimento facial estar consolidado em produ\c{c}\~ao, dois desafios continuam
relevantes para o aluno de engenharia: (i) entender o \emph{pipeline} completo,
da c\^amera ao classificador, em vez de tratar o reconhecimento como uma ``caixa-preta''; e
(ii) lidar com ataques de \emph{spoofing} (foto, v\'ideo, m\'ascara) usando t\'ecnicas de
\emph{prova de vida} (\emph{liveness detection}). Sistemas comerciais costumam combinar
v\'arias heur\'isticas geom\'etricas e modelos espec\'ificos para essa tarefa, e essa combina\c{c}\~ao
\'e justamente o que torna o problema interessante do ponto de vista did\'atico.

Este projeto prop\~oe a constru\c{c}\~ao de uma aplica\c{c}\~ao web completa, o \textbf{FaceID Lab},
que implementa o ciclo completo: \emph{captura} pela webcam, \emph{detec\c{c}\~ao} de faces,
\emph{extra\c{c}\~ao} de descritores de 128 dimens\~oes, \emph{verifica\c{c}\~ao}
por dist\^ancia euclidiana contra um banco cadastrado e \emph{prova de vida} para mitigar
ataques de \emph{spoofing}. A relev\^ancia social \'e direta: o mesmo n\'ucleo serve para
controle de presen\c{c}a em aulas, identifica\c{c}\~ao de pacientes em cl\'inicas e cadastro
facial em servi\c{c}os p\'ublicos -- todos casos em que aceitar uma foto impressa no lugar
de uma pessoa real seria uma falha grave.

% ==========================================================================
\section{Objetivos}

\subsection*{Objetivo Geral}

Desenvolver, avaliar e documentar um sistema web de reconhecimento facial \emph{end-to-end},
com cadastro guiado por prova de vida e identifica\c{c}\~ao por compara\c{c}\~ao de descritores
profundos, aplicando os conceitos de PDI ministrados na disciplina.

\subsection*{Objetivos Espec\'ificos}

\begin{enumerate}[leftmargin=*]
    \item Implementar a etapa de \textbf{aquisi\c{c}\~ao} via webcam no navegador, com
    pr\'e-processamento (redimensionamento e normaliza\c{c}\~ao) compat\'ivel com os modelos
    utilizados.

    \item Aplicar \textbf{detec\c{c}\~ao de face} usando o \emph{SSD MobileNet V1} no \emph{backend}
    e o \emph{Tiny Face Detector} no \emph{frontend}, comparando custo computacional e qualidade
    de detec\c{c}\~ao.

    \item Extrair \textbf{68 \emph{landmarks} faciais} e gerar o \textbf{descritor de 128 dimens\~oes}
    (\emph{face embedding}) por face detectada.

    \item Implementar a etapa de \textbf{decis\~ao por dist\^ancia} (dist\^ancia
    euclidiana entre descritores) com limiar configur\'avel, alinhada ao conte\'udo
    de m\'etodos de decis\~ao da disciplina.

    \item Implementar \textbf{prova de vida ativa} baseada em heur\'isticas geom\'etricas
    sobre os \emph{landmarks}: \emph{Eye Aspect Ratio} (EAR), \emph{yaw} estimado por
    rela\c{c}\~ao nariz-mand\'ibula e an\'alise de consist\^encia temporal dos descritores
    entre frames.

    \item Implementar \textbf{verifica\c{c}\~ao anti-adultera\c{c}\~ao} no servidor:
    todo descritor enviado pelo cliente \'e re-extra\'ido a partir da imagem original e
    confrontado com o que o cliente reportou, impedindo que um atacante envie um
    descritor manipulado sem refletir a face real.

    \item Implementar \textbf{detec\c{c}\~ao de cadastros duplicados}: antes de gravar,
    o novo descritor \'e comparado com todos os j\'a cadastrados; se a dist\^ancia
    for inferior a um limiar tolerante, o cadastro \'e recusado.

    \item Definir uma base de testes pr\'opria do autor (com consentimento)
    para \textbf{avaliar} taxa de falsa aceita\c{c}\~ao (FAR), taxa de falsa rejei\c{c}\~ao
    (FRR) e tempo m\'edio de identifica\c{c}\~ao, calibrando o limiar de decis\~ao.

    \item Disponibilizar uma \textbf{interface web utiliz\'avel} (cadastro com
    prova de vida, cadastro simples \emph{legado}, identifica\c{c}\~ao, listagem
    de pessoas e consulta individual por ID) e suportar o upload opcional de
    \textbf{imagem de documento} de identidade no cadastro, para \emph{auditoria
    posterior}.

    \item \textbf{Documentar} o pipeline completo no relat\'orio final.
\end{enumerate}

\subsection*{Resultados Esperados}

\begin{itemize}[leftmargin=*]
    \item Aplica\c{c}\~ao funcional rodando localmente (\emph{frontend} React + \emph{backend}
    Node.js), com tempo de identifica\c{c}\~ao inferior a 1\,s em CPU de \emph{notebook}.
    \item Taxa de acerto $\geq 95\%$ na base interna de testes para o limiar calibrado.
    \item Rejei\c{c}\~ao consistente de ataques triviais de \emph{spoofing} (foto impressa,
    foto em tela de celular est\'atica) atrav\'es da prova de vida.
    \item Relat\'orio final descrevendo cada bloco do \emph{pipeline} com os fundamentos
    te\'oricos correspondentes (\emph{e.g.} filtragem, normaliza\c{c}\~ao, classificador por
    dist\^ancia, m\'etodos estat\'isticos).
    \item Apresenta\c{c}\~ao com demonstra\c{c}\~ao ao vivo do cadastro com prova de vida
    e da identifica\c{c}\~ao em tempo real.
\end{itemize}

% ==========================================================================
\section{M\'etodos}

\subsection{Estado da Arte e Refer\^encias}

O reconhecimento facial moderno segue, em geral, tr\^es paradigmas:

\begin{enumerate}[leftmargin=*]
    \item \textbf{Eigenfaces / Fisherfaces} (\emph{Turk \& Pentland}, 1991; \emph{Belhumeur et al.},
    1997): proje\c{c}\~ao em subespa\c{c}o (PCA/LDA) das imagens da face. Simples, mas
    sens\'ivel a varia\c{c}\~oes de pose e ilumina\c{c}\~ao.

    \item \textbf{Descritores locais} (LBP, HOG): caracter\'isticas \emph{handcrafted} mais
    robustas a varia\c{c}\~oes locais, combinadas com classificadores cl\'assicos.

    \item \textbf{Descritores aprendidos por redes convolucionais profundas}:
    \emph{FaceNet} (\emph{Schroff et al.}, 2015), \emph{ArcFace} (\emph{Deng et al.}, 2019)
    e variantes, que treinam uma rede para mapear faces em um espa\c{c}o de
    \emph{embeddings} onde dist\^ancia $\equiv$ ``mesma pessoa''. \'E hoje o
    padr\~ao de fato.
\end{enumerate}

Para \emph{liveness detection} h\'a duas grandes linhas: (i) \emph{passive liveness}
(an\'alise de textura, refl ex\~oes especulares, modelos como \emph{Silent Face
Anti-Spoofing} da MiniVision); e (ii) \emph{active liveness} (pedir intera\c{c}\~oes
como piscar, virar a cabe\c{c}a, sorrir, e validar pela varia\c{c}\~ao temporal).

\subsection{Metodologia Adotada}

Adotamos o paradigma de \textbf{descritores profundos} via biblioteca
\texttt{face-api.js}/\texttt{@vladmandic/face-api}, que disp\~oe de redes pr\'e-treinadas
em JavaScript/TypeScript, executando tanto no \emph{frontend} (WebGL) quanto no
\emph{backend} (Node.js + \texttt{canvas}). Os modelos utilizados ser\~ao:

\begin{itemize}[leftmargin=*]
    \item \textbf{SSD MobileNet V1} (\emph{backend}) e \textbf{Tiny Face Detector}
    (\emph{frontend}) para detec\c{c}\~ao da face na imagem;

    \item \textbf{Face Landmark 68-net} para a localiza\c{c}\~ao dos 68 pontos faciais
    (olhos, sobrancelhas, nariz, boca, mand\'ibula);

    \item \textbf{FaceNet-like Recognition Net} para gerar o descritor de 128 dimens\~oes.
\end{itemize}

A \textbf{classifica\c{c}\~ao} segue o m\'etodo de decis\~ao por dist\^ancia visto na disciplina:
para cada pessoa cadastrada s\~ao armazenados \emph{m\'ultiplos} descritores
$\{\mathbf{c}_{i,j}\}_j$ (um por \emph{frame} aceito na captura). Para uma consulta com
descritor $\mathbf{q}$, calcula-se a dist\^ancia euclidiana
\[
    d_{i,j} = \| \mathbf{q} - \mathbf{c}_{i,j} \|_2 ,
\]
e identifica-se a pessoa $i^{*} = \arg\min_{i,j} d_{i,j}$ se $\min_{i,j} d_{i,j} < \tau$,
onde $\tau$ \'e o limiar calibrado experimentalmente. O valor inicial adotado
no arquivo de configura\c{c}\~ao do servidor \'e $\tau = 0{,}55$, e ser\'a refinado
empiricamente durante a fase de avalia\c{c}\~ao.

\subsection{Pipeline de Prova de Vida}

O cadastro do usu\'ario \'e feito atrav\'es de um m\'odulo \texttt{LivenessCapture} que
roda no navegador e coleta v\'arios \emph{frames} durante uma janela fixa de 5\,s.
A detec\c{c}\~ao em tempo real usa o \emph{Tiny Face Detector} (\texttt{inputSize}\,$=$\,416,
\texttt{scoreThreshold}\,$=$\,0,5), escolhido pelo baixo custo computacional adequado
\`a execu\c{c}\~ao em WebGL no \emph{frontend}. Para cada \emph{frame} calculamos:

\begin{itemize}[leftmargin=*]
    \item \textbf{EAR} (\emph{Eye Aspect Ratio}, Soukupov\'a \& \v{C}ech, 2016):
    \[
        \mathrm{EAR} = \frac{\|p_2 - p_6\| + \|p_3 - p_5\|}{2\, \|p_1 - p_4\|}
    \]
    aplicado aos \emph{landmarks} dos olhos, para detectar piscadas.

    \item \textbf{Estimativa de \emph{yaw}} a partir da raz\~ao horizontal entre a posi\c{c}\~ao
    do nariz e os extremos da mand\'ibula -- detecta rosto frontal vs.\ rota\c{c}\~ao
    lateral.

    \item \textbf{Consist\^encia dos descritores}: dist\^ancia m\'axima e m\'edia entre os
    descritores dos frames coletados. Frames id\^enticos pixel-a-pixel ou rostos diferentes
    s\~ao rejeitados.
\end{itemize}

S\~ao exigidos crit\'erios m\'inimos para aprovar a captura:
$\geq 15$ \emph{frames} v\'alidos; fra\c{c}\~ao $\geq 60\%$ de \emph{frames} com face presente;
faixa de \emph{yaw} compat\'ivel com micromovimento natural (entre $\sim$\,1,5\textdegree\ e
$\sim$\,18\textdegree, n\~ao zero -- para descartar fotos est\'aticas -- mas n\~ao excessiva);
dist\^ancia m\'axima inter-\emph{frame} dentro de \texttt{SAME\_PERSON\_MAX\_DIST}\,$=$\,0,5;
e dist\^ancia \emph{m\'edia} acima de um piso m\'inimo ($>$\,0,005), o que rejeita capturas
em que os frames s\~ao praticamente id\^enticos pixel-a-pixel (sintoma t\'ipico de
foto retransmitida ou v\'ideo congelado). Tr\^es \emph{frames} representativos
(in\'icio, meio e fim da janela) t\^em seus descritores persistidos para a
classifica\c{c}\~ao multi-amostra, e o \emph{frame} inicial \'e mantido como imagem
principal de refer\^encia da pessoa cadastrada.

\subsection{Defesas Adicionais no Servidor}

Al\'em das verifica\c{c}\~oes no cliente, o servidor implementa duas barreiras
complementares no fluxo de cadastro com prova de vida:

\begin{itemize}[leftmargin=*]
    \item \textbf{Anti-adultera\c{c}\~ao cliente-servidor:} para cada \emph{frame} recebido,
    o servidor re-extrai o descritor de 128-D a partir da imagem original e
    compara com o descritor que o cliente reportou. Se a dist\^ancia
    ultrapassar \texttt{MAX\_CLIENT\_SERVER\_DESCRIPTOR\_DIST}\,$=$\,0,45 o cadastro \'e
    rejeitado, evitando que um cliente malicioso fa\c{c}a o \emph{bypass} da prova de
    vida enviando descritores forjados.

    \item \textbf{Detec\c{c}\~ao de cadastro duplicado:} antes de gravar, todos os
    descritores rec\'em-extra\'idos s\~ao comparados com a base existente. Se a
    menor dist\^ancia for inferior a \texttt{DUPLICATE\_THRESHOLD}\,$=$\,0,5 o sistema
    reconhece a pessoa como j\'a cadastrada e rejeita o novo registro --
    impedindo que o mesmo usu\'ario abra m\'ultiplas identidades.
\end{itemize}

Como \textbf{evolu\c{c}\~ao planejada}, est\'a previsto integrar o modelo
\emph{Silent-Face-Anti-Spoofing} (MiniVision) convertido para ONNX e executado via
\texttt{onnxruntime-web} -- o c\'odigo j\'a possui um \emph{stub} de integra\c{c}\~ao
(\texttt{antispoof.ts}) preparado para receber esse classificador adicional.

\subsection{Etapas de PDI Cobertas}

A tabela a seguir relaciona cada bloco do projeto aos t\'opicos do conte\'udo program\'atico
da disciplina:

\begin{center}
\begin{tabular}{lp{8.5cm}}
\toprule
\textbf{T\'opico do plano de ensino} & \textbf{Onde \'e exercitado no projeto} \\
\midrule
Aquisi\c{c}\~ao e amostragem & Captura via \texttt{getUserMedia}, escolha de resolu\c{c}\~ao 640$\times$480 \\
Processamento pontual e filtragem & Normaliza\c{c}\~ao e redimensionamento antes do detector \\
Segmenta\c{c}\~ao / detec\c{c}\~ao & SSD MobileNet e Tiny Face Detector localizam a face \\
Detec\c{c}\~ao de bordas e descontinuidades & Base te\'orica para os 68 \emph{landmarks} \\
Descritores de caracter\'isticas & \emph{Embedding} de 128-D (FaceNet-like) \\
Reconhecimento por dist\^ancia & Classifica\c{c}\~ao 1-NN por dist\^ancia euclidiana sobre m\'ultiplos descritores por pessoa, com limiar configur\'avel \\
M\'etodos estat\'isticos / supervis\~ao & An\'alise FAR/FRR; modelos pr\'e-treinados com aprendizagem supervisionada \\
\bottomrule
\end{tabular}
\end{center}

% ==========================================================================
\section{Materiais e Ferramentas}

\subsection{Fontes de Imagens e Base de Dados}

\begin{itemize}[leftmargin=*]
    \item \textbf{Aquisi\c{c}\~ao em tempo real:} webcam do \emph{notebook} do autor.
    Resolu\c{c}\~ao alvo de 640$\times$480, $\sim$30\,fps.

    \item \textbf{Base de cadastros internos:} o autor e um pequeno grupo de
    volunt\'arios identificados ser\~ao cadastrados (com consentimento expl\'icito).
    Em cada cadastro, o m\'odulo de prova de vida coleta v\'arios \emph{frames}
    durante uma janela de 5\,s com o rosto frontal e seleciona automaticamente
    \textbf{3 \emph{frames} representativos} (in\'icio, meio e fim da janela), gerando
    portanto \textbf{3 \emph{embeddings} por pessoa} -- todos da pose frontal, mas com
    pequena varia\c{c}\~ao temporal de \emph{yaw} e ilumina\c{c}\~ao para refletir
    a din\^amica natural do rosto.

    \item \textbf{Base de testes de \emph{spoofing}:} ser\~ao usadas fotos impressas e fotos
    exibidas em tela de celular do pr\'oprio autor para validar a prova de vida.

    \item \textbf{Imagem de documento (opcional):} o cadastro suporta o envio adicional
    de uma foto do documento de identidade do usu\'ario. Essa imagem n\~ao participa
    do reconhecimento -- \'e apenas arquivada para auditoria humana posterior, simulando
    o requisito comum em cen\'arios reais (controle de acesso institucional,
    abertura de contas etc.).

    \item \textbf{Base p\'ublica auxiliar:} \emph{Labeled Faces in the Wild} (LFW) ou subconjunto
    do \emph{CelebA}, exclusivamente para refer\^encia comparativa de FAR/FRR, sem persistir
    cadastros.
\end{itemize}

\subsection{Bibliotecas, \emph{Frameworks} e Recursos}

\begin{center}
\begin{tabular}{llp{6cm}}
\toprule
\textbf{Componente} & \textbf{Ferramenta} & \textbf{Estado} \\
\midrule
Reconhecimento facial & \texttt{@vladmandic/face-api} (FaceNet-like) & Instalado \\
Detector pesado & SSD MobileNet V1 (backend) & Instalado \\
Detector leve & Tiny Face Detector (frontend) & Instalado \\
\emph{Landmarks} & Face Landmark 68-net & Instalado \\
Anti-spoof passivo & Silent-Face-Anti-Spoofing $+$ \texttt{onnxruntime-web} & Em planejamento \\
Frontend & React 18, TypeScript, Vite & Instalado \\
Backend & Node.js, Express, TypeScript & Instalado \\
Persist\^encia & Prisma ORM + SQLite & Instalado \\
Manipula\c{c}\~ao de imagens (backend) & \texttt{canvas}, \texttt{multer} & Instalado \\
Captura no navegador & \texttt{getUserMedia}, \texttt{<canvas>} & Instalado \\
Empacotamento opcional & Docker, \texttt{docker-compose} & Configurado \\
\bottomrule
\end{tabular}
\end{center}

\subsection{Hardware}

Nenhum hardware espec\'ifico precisa ser adquirido. O \emph{notebook} do autor,
com webcam integrada e CPU $\geq$ 4 \emph{cores}, \'e suficiente para a execu\c{c}\~ao em
tempo real (n\~ao h\'a depend\^encia de GPU dedicada).

\subsection{Licen\c{c}as e \'Etica}

Todo o \emph{software} utilizado \'e livre (MIT/Apache) ou j\'a est\'a no acervo dispon\'ivel
para discentes. O uso de imagens de faces respeita: (i) consentimento expl\'icito do autor e
dos volunt\'arios envolvidos; (ii) restri\c{c}\~ao da base de cadastros ao autor e aos
volunt\'arios identificados; (iii) op\c{c}\~ao de remo\c{c}\~ao a qualquer momento; e (iv)
n\~ao publica\c{c}\~ao de imagens ou descritores em reposit\'orios p\'ublicos.

% ==========================================================================
\section{Cronograma}

O cronograma a seguir foi alinhado ao calend\'ario oficial da disciplina EEL7815 em
2026.1 (in\'icio em 12/03, encerramento em 09/07).

\begin{center}
\begin{tabular}{cllp{6.8cm}}
\toprule
\textbf{Semana} & \textbf{Per\'iodo} & \textbf{Etapa} & \textbf{Entreg\'aveis} \\
\midrule
1     & 12/03         & Setup           & Leitura do plano de ensino; revis\~ao bibliogr\'afica inicial; defini\c{c}\~ao do escopo. \\
2--3  & 19/03--01/04  & Fundamenta\c{c}\~ao  & Estudo de detectores, \emph{landmarks} e \emph{embeddings}; instala\c{c}\~ao do ambiente; valida\c{c}\~ao do \emph{boilerplate} \texttt{face-api.js}. \\
4--5  & 02/04--15/04  & Backend         & API REST (Express + Prisma), endpoint de cadastro e identifica\c{c}\~ao, persist\^encia do descritor. \\
6--7  & 16/04--29/04  & Frontend b\'asico    & Telas de cadastro/identifica\c{c}\~ao em React, captura via webcam, integra\c{c}\~ao com a API. \\
8--9  & 30/04--13/05  & Liveness ativo  & Implementa\c{c}\~ao do \texttt{LivenessCapture} (EAR, yaw, consist\^encia de descriptors). \\
10--11 & 14/05--27/05 & Testes internos & Cadastro dos integrantes; medi\c{c}\~ao de FAR/FRR; calibra\c{c}\~ao de $\tau$. \\
12    & 28/05--03/06  & Anti-spoof (opcional) & Integra\c{c}\~ao do modelo Silent-Face-Anti-Spoofing via ONNX, se cronograma permitir. \\
13--14 & 04/06--17/06 & Avalia\c{c}\~ao final  & Bateria de testes de \emph{spoofing} (foto, tela), comparativos quantitativos. \\
15    & 18/06--24/06  & Relat\'orio    & Reda\c{c}\~ao do relat\'orio final (20\% da nota). \\
16    & 25/06         & Apresenta\c{c}\~ao    & Ensaio + apresenta\c{c}\~ao do projeto (80\% da nota). \\
17    & 02/07--09/07  & Reserva       & Eventuais ajustes, atividades de recupera\c{c}\~ao/substitutivas. \\
\bottomrule
\end{tabular}
\end{center}

\textit{Observa\c{c}\~ao:} as datas de apresenta\c{c}\~ao formal ser\~ao confirmadas com o professor
no decorrer do semestre, conforme o cronograma de aulas presenciais.

% ==========================================================================
\section*{Refer\^encias B\'asicas}

\begin{itemize}[leftmargin=*]
    \item \textbf{[1]} R.\,C. Gonzalez, R.\,E. Woods. \emph{Digital Image Processing}.
    Addison-Wesley Longman.

    \item \textbf{[2]} F. Schroff, D. Kalenichenko, J. Philbin. \emph{FaceNet: A Unified
    Embedding for Face Recognition and Clustering}. CVPR, 2015.

    \item \textbf{[3]} T. Soukupov\'a, J. \v{C}ech. \emph{Real-Time Eye Blink Detection using
    Facial Landmarks}. 21st Computer Vision Winter Workshop, 2016.

    \item \textbf{[4]} J. Deng \emph{et al.} \emph{ArcFace: Additive Angular Margin Loss for
    Deep Face Recognition}. CVPR, 2019.

    \item \textbf{[5]} S\'itio complementar: \emph{Image Processing Place}.

    \item \textbf{[6]} V. Mandic. \texttt{@vladmandic/face-api}. Documenta\c{c}\~ao oficial.
\end{itemize}

\end{document}
